% from ega0-1-fr.tex
\oldpage[0\textsubscript{I}]{75}

ces homomorphismes forment un système inductif~; par passage à la limite
inductive, ils définissent donc un homomorphisme canonique
$A_{\{S\}}\to A\{S^{-1}\}$.
\end{env}

\begin{proposition}[7.6.16]
\label{0.7.6.16-fr}
Si $A$ est un anneau noethérien, $A\{S^{-1}\}$ est un module plat sur
$A_{\{S\}}$.
\end{proposition}

En effet (7.6.14), $A\{S^{-1}\}$ est plat sur chacun des anneaux
$A_{\{f\}}$ pour $f\in S$, et la conclusion résulte de (6.2.3).

\begin{proposition}[7.6.17]
\label{0.7.6.17-fr}
Soit $\mathfrak p$ un idéal premier ouvert dans un anneau admissible $A$, et
soit $S=A-\mathfrak p$. Alors les anneaux $A\{S^{-1}\}$ et $A_{\{S\}}$ sont
des anneaux locaux, l'homomorphisme canonique
$A_{\{S\}}\to A\{S^{-1}\}$ est local et les corps résiduels de $A_{\{S\}}$
et $A\{S^{-1}\}$ sont canoniquement isomorphes au corps des fractions de
$A/\mathfrak p$.
\end{proposition}

En effet, soit $\mathfrak J\subset\mathfrak p$ un idéal de définition de
$A$~; on a $S^{-1}\mathfrak J\subset S^{-1}\mathfrak p=\mathfrak p
A_{\mathfrak p}$, donc $A_{\mathfrak p}/S^{-1}\mathfrak J$ est un anneau
local~; on conclut de (7.1.12), (7.6.9) et (7.6.11, (i)) que
$A\{S^{-1}\}$ est un anneau local. Posons
$\mathfrak m=\varinjlim_{f\in S}\mathfrak p_{\{f\}}$, qui est un idéal de
$A_{\{S\}}$~; nous allons voir que tout élément de $A_{\{S\}}$ n'appartenant
pas à $\mathfrak m$ est inversible. En effet, un tel élément est l'image dans
$A_{\{S\}}$ d'un élément $z\in A_{\{f\}}$ n'appartenant pas à
$\mathfrak p_{\{f\}}$, pour un $f\in S$~; son image canonique $z_0$ dans
$A_{\{f\}}/\mathfrak J_{\{f\}}=S_f^{-1}(A/\mathfrak J)$ n'appartient donc
pas à $S_f^{-1}(\mathfrak p/\mathfrak J)$ (7.6.9), ce qui signifie que
$z_0=\overline{x}/\overline{f}^k$, où $x\notin\mathfrak p$ et
$\overline{x},\overline{f}$ sont les classes de $x,f$ mod.~$\mathfrak J$.
Comme $x\in S$, on a $g=xf\in S$, et dans $S_g^{-1}A$,
$y_0=x^{k+1}/g^k$, image canonique de $x/f^k\in S_f^{-1}A$, admet un
inverse $x^{k-1}f^{2k}/g^k$. Cela entraîne \emph{a fortiori} que l'image de
$y_0$ dans $S_g^{-1}A/S_g^{-1}\mathfrak J$ est inversible, donc (7.6.9 et
7.1.12) l'image canonique $y$ de $z$ dans $A_{\{g\}}$ est inversible~;
l'image de $z$ dans $A_{\{S\}}$ (égale à celle de $y$) est par suite
inversible. On voit donc que $A_{\{S\}}$ est un anneau local d'idéal maximal
$\mathfrak m$~; en outre, l'image de $\mathfrak p_{\{f\}}$ dans
$A\{S^{-1}\}$ est contenue dans l'idéal maximal $\mathfrak p\{S^{-1}\}$ de
cet anneau~; \emph{a fortiori}, l'image de $\mathfrak m$ dans
$A\{S^{-1}\}$ est contenue dans $\mathfrak p\{S^{-1}\}$, donc
l'homomorphisme canonique $A_{\{S\}}\to A\{S^{-1}\}$ est local. Enfin,
comme tout élément de $A\{S^{-1}\}/\mathfrak p\{S^{-1}\}$ est image d'un
élément d'un anneau $S_f^{-1}A$ pour un $f\in S$ convenable,
l'homomorphisme $A_{\{S\}}\to A\{S^{-1}\}/\mathfrak p\{S^{-1}\}$ est
surjectif, et donne donc par passage aux quotients un isomorphisme des corps
résiduels.

\begin{corollary}[7.6.18]
\label{0.7.6.18-fr}
Sous les hypothèses de (7.6.17), si on suppose de plus que $A$ est un anneau
adique noethérien, les anneaux locaux $A\{S^{-1}\}$ et $A_{\{S\}}$ sont
noethériens, et $A\{S^{-1}\}$ est un $A_{\{S\}}$-module fidèlement plat.
\end{corollary}

On sait déjà (7.6.11, (ii)) que $A\{S^{-1}\}$ est noethérien et
$A_{\{S\}}$-plat (7.6.16)~; comme l'homomorphisme
$A_{\{S\}}\to A\{S^{-1}\}$ est local, on en conclut que $A\{S^{-1}\}$ est
un $A_{\{S\}}$-module fidèlement plat (6.6.2), et par suite que $A_{\{S\}}$
est noethérien (6.5.2).

\subsection{Produits tensoriels complétés}
\label{subsection:0.7.7-fr}

\begin{env}[7.7.1]
\label{0.7.7.1-fr}
Soient $A$ un anneau linéairement topologisé, $M$, $N$, deux $A$-modules
linéairement topologisés. Soient $\mathfrak J$, $V$, $W$ des voisinages
ouverts de $0$ dans $A$, $M$, $N$ respectivement, qui soient des
$A$-modules, et tels que $\mathfrak JM\subset V$, $\mathfrak JN\subset W$,
de sorte que $M/V$ et $N/W$ peuvent être considérés comme des
$(A/\mathfrak J)$-modules. Lorsque $\mathfrak J$, $V$, $W$ parcourent les
systèmes de voisinages ouverts vérifiant les conditions précédentes, il est
immédiat que les modules
$(M/V)\otimes_{A/\mathfrak J}(N/W)$ forment un système projectif de modules

